{\centering\tiny\renewcommand{\arraystretch}{0.92}
\begin{tabularx}{\textwidth}{@{}>{\raggedright\arraybackslash}p{2.6cm}Y@{}}
\toprule
family & rules (id: name [status]) \\
\midrule
\textsc{sv} (15) svara (vowel) yati & \textbf{01} vowel-class yati [C]; \textbf{01.1} consonant AND vowel must pair [M]; \textbf{02} svara-pradhāna (sandhi) yati [C]; \textbf{02.2} drutam+V and yaḍāgama+V read as the vowel [C]; \textbf{02.10} vowel-only acceptance [R]; \textbf{03.1} bare \te{ఋ} with I-class vowels [C]; \textbf{03.2} \te{ఋ} as second-word initial [C]; \textbf{04} hidden-\te{అ} (visarga-ō) yati [C]; \textbf{04.4} hidden \te{అ} by pūrvarūpa [I]; \textbf{05} \te{ఋ} with \te{రి రీ రె రే} [C]; \textbf{05.2} bound \te{ఋ} with \te{రి}-series [C]; \textbf{05.5} \te{ఋ} with \te{ఱి}-series [O]; \textbf{06} bound \te{ఋ} with a bare I-class vowel [C]; \textbf{07} bound \te{ఋ} with bound \te{ఋ} [C]; \textbf{08} vṛddhi dual option [C] \\
\textsc{vy} (16) vyañjana (consonant) yati & \textbf{01} identity yati [C]; \textbf{02.1} same-varga stop yati [C]; \textbf{03} anusvāra+stop with class nasal [C]; \textbf{04} na-ṇa yati [C]; \textbf{05} anusvāra+ṭa/ta stop with the other nasal [C]; \textbf{06} anusvāra+ṭa with anusvāra+ta [C]; \textbf{08} ya-ha yati [C]; \textbf{09} vowel with ya [C]; \textbf{10} vowel with ha [C]; \textbf{11} sibilant yati [C]; \textbf{12} sibilant with palatal [C]; \textbf{13} va-ba yati [C]; \textbf{14} va with pa pha bha [C]; \textbf{15} la-ḷa yati [C]; \textbf{16} mu with pu-series [C]; \textbf{17} stem mu with pu-series [C] \\
\textsc{sy} (7) saṃyukta (conjunct) & \textbf{01} any cluster constituent [C]; \textbf{03} geminate read as its consonant [C]; \textbf{05} kṣa split [C]; \textbf{08} drutam fusion across the line break [C]; \textbf{09} pollu-la fusion [C]; \textbf{11} same constituent for every caesura [M]; \textbf{11.2} sīsamu halves exempt [C] \\
\textsc{sp} (6) special consonants & \textbf{01} jña with na/ṇa [C]; \textbf{02} jña with ka-varga [C]; \textbf{03} anusvāra+semivowel/spirant with ma [C]; \textbf{06} elided \te{అది/అవి} [M]; \textbf{07} ra only with ra, \b{r}a only with \b{r}a [C]; \textbf{08} ra with \b{r}a [O] \\
\textsc{ub} (16) ubhaya: a vowel and a consonant track & \textbf{01} -iñcu dual track [C]; \textbf{02} -ēni dual track [C]; \textbf{03} -ēsi dual track [C]; \textbf{04} -eḍu measure dual track [C]; \textbf{05} ordinal -ava dual track [C]; \textbf{10} upasarga dual track [C]; \textbf{11.1} nitya-samāsa dual track [C]; \textbf{11.3} sa- prefix dual track [M]; \textbf{11.4} viśvāmitra vowel track (flagged) [R]; \textbf{12} deity-name dual track [C]; \textbf{13} nañ-compound dual track [C]; \textbf{14} kyac vowel track forbidden [F]; \textbf{16} native fused-unit dual track [C]; \textbf{17} name+honorific dual track [C]; \textbf{18} rugāgama dual track [C]; \textbf{19} dative kai dual track [C] \\
\textsc{py} (4) prāsa-yati & \textbf{01} prāsa-yati fallback [C]; \textbf{02} permitted metres only [M]; \textbf{03} weight parity of the preceding aksharas [M]; \textbf{25} head yati takes precedence [C] \\
\textsc{dl} (10) derived and later rules & \textbf{01} ḷ-vowel yati [C]; \textbf{02} ra-la yati [R]; \textbf{02.1} ra-ḷa yati [R]; \textbf{04} la/ḷa with ḍa [C]; \textbf{05} da-ḍa yati [R]; \textbf{06} mu with vu [C]; \textbf{09} snā with ta [R]; \textbf{10} augment consonant never on the hal track [M]; \textbf{11} geminate ṭa [M]; \textbf{12} verbal -eḍu [M] \\
\textsc{rj} (12) rejected pairs & \textbf{00} no maitri between these consonants [F]; \textbf{03} vowel cannot leave its consonant (except \te{ఋ}) [M]; \textbf{04} vowel classes differ [M]; \textbf{05} stop with class nasal without anusvāra [F]; \textbf{06} semivowel/spirant cross pairs [F]; \textbf{07} ma with va [F]; \textbf{08} u-vowel with vu [F]; \textbf{10} ṛ with ru [F]; \textbf{12} ṛ with li [F]; \textbf{13} la/ḷa with ḍha [F]; \textbf{14} other dental-retroflex pairs [F]; \textbf{25} constituent switched between caesuras [D] \\
\textsc{ep} (6) engine procedure & \textbf{02} normalisation [M]; \textbf{03} candidate lattice [C]; \textbf{03.1} blocked readings [M]; \textbf{04} compare as (consonant class, vowel class) [C]; \textbf{07} bahu-yati check [M]; \textbf{08} prāsa-yati fallback [C] \\
\bottomrule
\end{tabularx}
\captionof{table}{The yati rules (\texttt{yati\_rules.yaml}), distilled from chapter 6 of Padya Vidya \citep{padyavidya} (92 rules). Status: C canonical, Cs canonical subtype, M mandatory (all profiles), R accepted relaxation (relaxed and historical profiles), O orthographic relaxation, D deprecated (historical only), F forbidden (a named failure), I informational.}
\label{tab:yatirules}\par}
\medskip
